% ---- pg-577 （书页 565） ----

\mvsub{最后评论}

对应分析本质上是一种图方法, 用于表示在低维空间中变量间的关系. 此方法可以看做是一种标度方法, 因此可以看做是其他方法如多维标度法(12.6 节)或双重信息法(12.8 节)的补充. 对应分析与主成分分析(第 10 章)和正则相关分析(第 10 章)有联系. 欲了解更多有关对应分析的知识, 可阅读格里纳克的书(见文献[14]).

\section{用于观察抽样单元和变量的双重信息图}
\phantomsection\label{sec:12.8}

双重信息图是 $n\times p$ 数据集中信息的一种图形表示方法. “双重”是指包含在数据集中的两类信息: 行中的信息属于样本或抽样单元, 而列中的信息则属于变量.

当只有两个变量时, 通过散布图就能将抽样单元中的信息和变量中的信息在同一张图上表示出来. 这样人们就可以通过视觉来检查一个抽样单元相对另一个抽样单元的位置, 以及两个变量中的每一个与任一个单元位置的相对重要性.

当存在若干个变量时, 我们可以构造散布图矩阵, 但不能将抽样单元表示在一张图上. 另一方面, 抽样单元的二维图可以通过对前两个主成分作图来得到(见 8.4 节). 隐藏在双重信息图背后的想法是将有关变量的信息加到主成分图中去.

\hyperref[fig:12.23]{图 12.23} 给出\hyperref[tab:12.4]{表 12.4} 中公用事业公司数据的双重信息图. 从图中可以看出公司之间是如何分组聚类的, 以及变量对它们位置的贡献. 例如, 变量 $X_4$(=年载荷因子)和 $X_8$(=总燃料成本)基本上将大部分沿海公司集中在图的右下角. 变量 $X_1$(=固定费用比)和 $X_2$(=资本回报率)将佛罗里达和路易斯安那的公司聚集在一起.

\begin{mvfigure}
\mvinc{pg-577}{0.229,0.560,0.786,0.876}
\phantomsection\label{fig:12.23}\mvcap{图 12.23\quad 公用事业公司数据的双重信息图}
\end{mvfigure}
